\documentclass[11pt]{article}
\usepackage[margin=2.2cm]{geometry}
\usepackage[most]{tcolorbox}
\usepackage{hyperref}
\usepackage{enumitem}
\usepackage{xcolor}

\tcbset{
  boxrule=1pt, arc=3pt,
  left=6pt, right=6pt, top=5pt, bottom=5pt,
  breakable,
  colback=white,
  fonttitle=\bfseries\normalsize,
}

\newtcolorbox{responsebox}{
  colback=blue!4,
  colframe=blue!70!black,
  colbacktitle=blue!18,
  coltitle=blue!50!black,
  title={Our response},
  before skip=0pt, after skip=0pt,
}

\newtcolorbox{reviewerbox}{
  colback=gray!3,
  colframe=gray!55,
  colbacktitle=gray!12,
  coltitle=gray!45!black,
  title={Reviewer comment},
  fontupper=\itshape,
  after skip=0pt,
}
\newcommand{\rsp}{\par\vspace{0.9em}}
\newcommand{\rc}[1]{\begin{reviewerbox}#1\end{reviewerbox}}
\newcommand{\rcitem}[2]{\begin{reviewerbox}\textbf{#1}\par\smallskip#2\end{reviewerbox}}
\newcommand{\rcbothq}[3]{%
  \begin{reviewerbox}\textbf{#1}\par\smallskip\textbf{Reviewer \#1:} #2\par\smallskip%
  \textbf{Reviewer \#2:} #3\end{reviewerbox}}
\makeatletter
\newcommand{\response}[2]{%
  \begin{responsebox}#1%
  \def\@responseitems{#2}%
  \ifx\@responseitems\empty\else\begin{itemize}[nosep,leftmargin=1.2em]#2\end{itemize}\fi
  \end{responsebox}\rsp}
\makeatother

\title{Point-by-Point Response to Reviewers\\[0.3em]
\large Manuscript PR-D-26-04476\\
\textit{Interpretable supervised feature extraction for time series and spatial data}}
\date{\today}

\begin{document}
\maketitle

\noindent Dear Editor, Associate Editor, and Reviewers,

We thank you for the constructive major-revision decision and for the detailed feedback, which has
substantially improved the manuscript. Below we respond to every point. We address
\textbf{Reviewer~\#2}'s suggestions through the editorial questionnaire (Q1--Q8) and answer
\textbf{Reviewer~\#1}'s metareview point by point. For each item, the lead sentence summarizes what
we changed; bullets quote key wording from the \textbf{revised} manuscript (in quotation marks).

The revised manuscript complies with the journal's \textbf{35-page limit}. Headline changes include: 
numbered Objectives and a condensed Related Work with expanded Table~1;
Methods renamed to \emph{Iterative Feature Discovery} (fit/selection halves); baseline
hyperparameters and dataset Table~2 in the main text; separate Supplementary Data (preprocessing
only); paired Wilcoxon tests and per-patient AZT1D reporting; multi-fold ablations with a
\emph{Resampling} variant; expanded REMC/AZT1D SHAP interpretation; 
and a general reframing regarding algorithm superiority and clinical applicability. 
Several requests could not be implemented at full depth in the main text without
exceeding that limit (e.g.\
step-by-step preprocessing for all datasets, multi-seed repeats of the full feature search,
95\% confidence-interval columns, additional spatial benchmarks, SHAP on every dataset, or
high-resolution supplementary figure files). Where we defer detail, we say so explicitly and provide
the closest alternative---usually the separate Supplementary Data document or a condensed main-text
summary.

\rcbothq{1. Are the objectives and the rationale of the study clearly stated?}{Yes.}
{The manuscript presents a clear overall motivation---addressing the trade-off between predictive
performance and interpretability in time series feature extraction---and this is reasonably well
justified in the introduction; however, the objectives are not explicitly stated and remain
implicit throughout the text.
To improve clarity, the authors should include a concise and explicit statement of the main
objective at the end of the introduction, more clearly separate the problem motivation from the
methodological contribution, and articulate specific research goals or guiding questions.
In addition, the novelty of the approach should be more sharply defined relative to existing
methods (e.g., interval-based or shapelet approaches), and the stated objectives should be
better aligned with the evaluation, particularly with respect to how interpretability is assessed.}

\response{We now state the study objectives explicitly. The revision adds Sec.~1.1
(Objectives) with three aims and a condensed Sec.~1.2 (Related Work) before Methods.}{
  \item Sec.~1.1 (new): ``The objectives of our study are as follows: (i) we present SublimeX \ldots{}
    (ii) we benchmark SublimeX against other existing methods on seven public biomedical datasets,
    and (iii) we assess robustness, generalization and interpretability through post-hoc analyses
    and ablations.''
  \item Sec.~1.2 (revised): ``SublimeX combines traceable segment means with supervised search over
    channel, transform, and location''; Table~1 compares methods by type, approximate feature count,
    multivariate handling (Yes = joint across channels; No = per-channel concatenation), and
    interpretability (High/Medium/Low with brief rationale).
}

\rcbothq{2. If applicable, is the application/theory/method/study reported in sufficient detail
to allow for its replicability and/or reproducibility?}{Yes [X].}
{The manuscript provides a generally solid methodological description of the proposed framework
(including the optimization procedure, transformations, and evaluation setup), which supports
partial reproducibility; however, several details remain insufficient for full replication.
To improve reproducibility, the authors should (1) specify all implementation details more
precisely, including exact hyperparameters of the downstream models (beyond default mentions) and
stopping criteria, (2) clarify how segment boundaries are sampled and constrained during
optimization, (3) provide full details of preprocessing steps for each dataset within the main
text rather than relying on supplementary material, (4) describe the exact data splits or make
them publicly available to ensure consistent benchmarking, (5) report random seeds and clarify
whether results are sensitive to initialization, (6) include computational cost details such as
total runtime per dataset and hardware variability, and (7) provide direct access to code,
configuration files, and scripts used to reproduce experiments, ideally with instructions for
end-to-end replication.}

\response{We expanded Methods so the study can be reproduced end to end. Because of the
  35-page limit, we kept the per-dataset preprocessing to a \textbf{separate Supplementary Data} file.
  However, we state all Baseline parameters clearly in the main text. 
  }{
  \item Methods, Iterative Feature Discovery (renamed from ``Segment-Based Feature Optimization''):
    ``center and width on a normalized $0$--$1$ scale'';
    a fixed 50/50 split of the training data into a ``\emph{fit half}, used to train the downstream
    model, and a \emph{selection half}, used to score candidate features'';
    ``terminates once no candidate yields a strict improvement''; fit/selection terminology replaces
    the earlier search-train/search-validation wording.
  \item Methods, Baselines: tsfresh, catch22, MiniRocket, RDST, and CNN hyperparameters in prose
    (replacing the former supplementary ``Baseline Implementations'' appendix).
  \item Methods, Implementation Details: ``300 trials per proposed feature''; ``maximum tree depth
    5''; ``published as an open-source Python package on PyPI (version~0.1.2)''; ``MacBook Pro M2
    (12 cores, 32\,GB RAM).''
  \item Table~3: now reports performance, feature counts, and wall-clock time.
  \item Manuscript repository: \url{https://github.com/Prgrmmrjns/SublimeX-manuscript} contains 
  step-by-step reproduction instructions (data download, experiment
    pipeline, tables/figures, LaTeX build)
}

\rcbothq{3. If applicable, are statistical analyses, controls, sampling mechanism, and statistical
reporting (e.g., P-values, CIs, effect sizes) appropriate and well described?}{Yes [X].}
{The manuscript includes standard evaluation practices (e.g., cross-validation, reporting mean
$\pm$ standard deviation across folds) and uses appropriate performance metrics depending on the
task (AUC, accuracy, RMSE), which supports a generally acceptable level of statistical reporting;
however, the statistical analysis lacks depth and rigor in several aspects.
To improve this, the authors should (1) include statistical significance testing when comparing
methods (e.g., paired tests across folds or datasets) to support claims of competitive
performance, (2) report confidence intervals in addition to standard deviations to better
quantify uncertainty, (3) clarify whether hyperparameter tuning and feature selection are nested
within cross-validation to avoid potential data leakage, (4) provide more detail on the sampling
strategy and class balance handling for each dataset, particularly in imbalanced settings, (5)
assess variability across runs by repeating experiments with different random seeds where
applicable, and (6) better justify the choice of evaluation protocols (e.g., single split
vs.\ cross-validation) and discuss their potential impact on the reported results.}

\response{We strengthened statistical reporting with paired Wilcoxon tests, explicit evaluation
protocols, and anti-leakage wording. For cross-validated benchmarks, tests pair matched folds;
for AZT1D we fully re-ran all methods so each of 23 patients contributes one held-out test RMSE
and significance is assessed on that patient-level distribution (not across folds). Table~3 reports
AZT1D as mean$\pm$SD over those 23 per-patient errors (e.g.\ SublimeX $30.98 \pm 5.65$ vs.\
MiniRocket $30.76 \pm 5.38$, $p{=}0.39$). We do not add separate 95\% confidence-interval columns
(five folds yield wide intervals, largely redundant with mean$\pm$SD) or repeat the full feature
search across random seeds.}{
  \item Methods, Evaluation: ``paired two-sided Wilcoxon tests ($\alpha{=}0.05$)'' on matched
    folds, LOSO subjects, or patients; Table~3 asterisks; anti-leakage wording; PAMAP2
    once-assigned fit/selection split.
  \item Methods, Evaluation (AZT1D): temporal 80/20 split per patient; pooled segment discovery
    on training windows only; RMSE as mean$\pm$SD over $n{=}23$ per-patient test errors (full
    re-run for all methods).
  \item Methods, Iterative Feature Discovery: fit/selection 50/50 split during search; test scoring
    in Evaluation.
  \item Table~4: multi-fold ablations with mean$\pm$SD (replacing single-fold ablations).
}

\rcbothq{4. Could the manuscript benefit from additional tables or figures, or from improving or
removing (some of the) existing ones?}{No.}
{The manuscript includes several relevant figures and tables that support the methodology and
results (e.g., workflow diagram, performance comparison, and feature analysis), but their
effectiveness could be improved.
Specifically, (1) the graphical abstract is visually appealing but somewhat generic and could be
made more informative by explicitly illustrating the optimization process, (2) Figure~1
(methodology) is helpful but visually dense and would benefit from simplification or clearer
labeling of key steps, (3) Figure~2 and Figure~3 provide useful insights but are difficult to
interpret due to small text and multiple panels, so readability should be improved (e.g., larger
fonts, clearer legends), (4) Table~3 would benefit from including statistical significance
indicators or highlighting meaningful differences beyond bolding best values, (5) an additional
table summarizing the number and type of extracted features per dataset would strengthen
interpretability claims, and (6) a more explicit visualization comparing interpretability (not
only performance) across methods would further support the main contribution.}

\response{We improved the graphical abstract and Figures~1--4 for readability and extended Tables~1 and~3 (Wilcoxon markers,
runtimes, and new Multivariate and Interpretability columns).}{
  \item Graphical abstract: redrawn with larger font size and SublimeX workflow from flowchart
  \item Figure~1: redrawn flowchart; caption states the loop stops when ``evaluation no longer
    improves.''
  \item Figure~2: regenerated with larger fonts (transform composition and pattern stability).
  \item Figure~3: incremental train/validation/test curves; AZT1D uses pooled training windows
    matching the feature-discovery protocol.
  \item Figure~4: full-width six-panel SHAP layout for REMC and AZT1D, only top 5 features instead of all features.
  \item Table~3: ``wall-clock time'' row and significance asterisks.
  \item Table~1: added Multivariate (Yes/No) and shortened Interpretability columns (High/Medium/Low
    with brief rationale); method types written out (Segment, Interval, Convolution, DL, etc.).
}

\rcbothq{5. Have the authors clearly emphasized the strengths of their
study/theory/methods/argument?}{Yes.}
{The manuscript demonstrates several clear strengths---particularly the balance between
interpretability and performance, the modular design of the framework, and the use of diverse
datasets---but these strengths are not always emphasized as clearly or strategically as they
could be.
To improve this, the authors should (1) explicitly summarize the main strengths of the method in
a dedicated contributions paragraph at the end of the introduction, (2) more clearly highlight the
novelty of combining supervised feature search with interpretable, segment-based representations,
(3) emphasize the flexibility and modularity of the framework as a practical advantage for
different domains, (4) reinforce the importance of compact feature sets as a key contribution
rather than a secondary outcome, and (5) more directly connect the empirical results and SHAP
analysis to the claimed strengths in interpretability and real-world applicability.}

\response{We make the key strengths---modularity and compact, traceable features---explicit in the
Objectives and Implementation Details, while keeping the claims balanced against the limitations.
The Discussion and Conclusion use measured, ``competitive'' wording.}{
  \item Sec.~1.1 (Objectives): compact traceable features linked to ablations and interpretability
    analyses.
  \item Table~1: multivariate Yes/No and interpretability ratings (High/Medium/Low) make the
    performance--interpretability trade-off explicit across method families.
  \item Methods, Implementation Details: ``modular, extensible design''; customizable objective,
    model, and transform library.
  \item Discussion and Conclusion:rephrased ``competitive'' wording; REMC before AZT1D in Summary;
    PAMAP2 LOSO interpretation in Discussion.
}

\rcbothq{6. Have the authors clearly stated the limitations of their
study/theory/methods/argument?}{Yes.}
{The manuscript acknowledges several important limitations---such as computational cost, potential
overfitting during feature search, and challenges in cross-subject generalization---which is a
clear strength; however, these limitations could be stated more explicitly and systematically.
To improve this, the authors should (1) consolidate the limitations into a clearly structured
paragraph or subsection for better visibility, (2) more explicitly discuss the impact of the
greedy search strategy and its susceptibility to local optima, (3) elaborate on the risk of
overfitting due to the fixed train/validation split used during optimization, (4) clarify the
scalability limitations in terms of runtime and high-dimensional search space, (5) more
critically address the limited comparison with additional state-of-the-art methods, and (6)
discuss the potential limitations of interpretability, particularly whether simple segment-based
features are sufficient for capturing complex patterns in all domains.}

\response{We expanded Discussion, Limitations: greedy search/local optima (with parallel ablation
caveat), PAMAP2 validation--test gap (${\sim}26\%$) and fixed-split risk supported by the
\emph{Resampling} ablation, omitted subject-holdout search for compute cost, slower runtime vs.\
lighter baselines (Table~3), and narrower TSC/UCR comparison scope. Future Directions adds nested
internal validation.}{
  \item Discussion, Limitations: greedy search; PAMAP2 gap; \emph{Resampling}; runtime; TSC/UCR scope.
  \item Discussion, Future Directions: nested or repeated internal validation.
}

\rcbothq{7. Does the manuscript structure, flow or writing need improving (e.g., the addition of
subheadings, shortening of text, reorganization of sections, or moving details from one section
to another)?}{No.}
{The manuscript is generally well organized and follows a standard structure, but the clarity and
flow can be improved in several areas.
Specifically, (1) the introduction would benefit from a clearer closing paragraph summarizing the
main contributions in a structured way, (2) some sections---particularly the
methodology---are dense and could be streamlined by reducing repetition and improving transitions
between subsections, (3) key implementation details currently placed in supplementary material
should be briefly summarized in the main text to improve readability, (4) the results section
could be better structured by explicitly linking each subsection to a specific research objective
or claim, and (5) the discussion could be made more concise by reducing redundancy and more
clearly separating interpretation, limitations, and future work.}

\response{We reorganized the manuscript for clearer flow within the 35-page limit: explicit
Objectives, a condensed Related Work, a separate Supplementary Data file, and a four-part
Discussion.}{
  \item Sec.~1.1--1.2: explicit objectives and a streamlined related-work contrast before Methods,
    with Table~1 summarizing method type, feature count, multivariate handling, and interpretability;
    smoother transitions between method families.
  \item Methods: ``Segment-Based Feature Optimization'' renamed \emph{Iterative Feature Discovery};
    reduced repetition and improved transitions across Signal Transformation, Iterative Feature
    Discovery, and Implementation Details.
  \item Methods, Datasets: Table~2 in main text (short caption; ``channels $\times$ length'');
    MIMIC-IV initial-feature block described; preprocessing in Supplementary Data.
  \item Methods, Baselines: library and CNN hyperparameters in main text (Section~Baselines);
    LightGBM readout in Implementation Details (depth~5).
  \item Results: Performance Comparison, Transform Composition, Incremental Feature Analysis
    (PAMAP2 numbers only; LOSO interpretation in Discussion), Ablation Study, Domain-Specific
    Feature Interpretation as distinct subsections.
  \item Discussion: Summary, Comparison to Related Work, Limitations, and Future Directions as
    separate subsections.
}

\rcbothq{8. Could the manuscript benefit from language editing?}{Yes.}{Yes.}

\response{We carried out a manuscript-wide language and flow edit: softened overstated wording,
standardized dataset names, and improved readability while staying within the 35-page limit (the
Abstract was condensed to $\leq$250 words).}{
  \item Introduction: tighter motivation and a clearer contribution lead-in (``To this end we
    present \textbf{SublimeX}\ldots''); Related Work now reads as a single contrast that ends with
    ``SublimeX bridges these strands by pairing traceable segment means with supervised search.''
  \item Methods: smoother Signal Transformation and Iterative Feature Discovery paragraphs and a
    crisper Implementation Details (modularity) passage, with no change to reported settings.
  \item Discussion: shorter paragraphs; opens with a SublimeX recap and measured performance
    language; separate Limitations / Future Directions subsections.
  \item Conclusion (revised): three generalized sentences on competitive traceable features,
    interpretability vs.\ performance, and open-source availability.
  \item Throughout: consistent dataset naming (MIMIC-IV, MIT-BIH, etc.) in Tables~2--4 and text.
}

\rcitem{Associate Editor (Reviewer \#1) --- full metareview}{The manuscript ``Interpretable
supervised feature extraction for time series and spatial data'', I would recommend major
revision. The topic is relevant and the proposed SublimeX framework has potential, but several
methodological, validation, and presentation issues need to be addressed before publication.

The manuscript presents an interesting interpretable feature extraction framework, but the novelty
of SublimeX should be stated more clearly. The authors should explain how it differs
substantially from existing interval-based, shapelet-based, ROCKET/MiniROCKET, and supervised
feature extraction methods.
The title mentions both ``time series and spatial data,'' but most of the paper focuses on
time-series biomedical datasets. The spatial data contribution is not sufficiently demonstrated.
The authors should either strengthen the spatial-data evaluation or revise the title and scope.
The abstract is informative, but it overstates competitiveness. SublimeX is not the best-performing
method on several datasets, including MITBIH, Emotions, MIMIC-IV, PAMAP2, and SVD. The abstract
should more cautiously describe the results.
The methodology needs more detail on the stopping criterion. The paper states that the loop stops
when no candidate improves performance, but the exact definition of improvement, tolerance
threshold, and validation metric direction should be clearly specified.
The use of a single deterministic 50/50 train/validation split during feature optimization may
introduce instability and overfitting. The authors should justify this choice and discuss whether
nested cross-validation or repeated validation would provide more reliable feature selection.
The internal feature-search validation split does not match the external evaluation design for
some datasets, especially PAMAP2 with leave-one-subject-out validation. This likely explains the
large validation-test gap. The authors should consider using subject-independent validation
during feature search.
The manuscript should clarify whether feature extraction is performed independently within each
training fold only. If any information from validation or test folds influences feature discovery,
this would create data leakage.
The comparison with baselines is incomplete. Important state-of-the-art time series methods such
as HIVE-COTEv2, InceptionTime, and LSTM/GRU-based models are only mentioned as limitations but
not evaluated. At minimum, the authors should justify why these were excluded.
The authors should provide statistical significance testing across methods. Many reported
differences are very small, such as AZT1D RMSE 31.11 vs.\ 31.18 and REMC AUC 0.931 vs.\ 0.931, so
claims of superior performance require statistical support.
The results show that SublimeX is competitive but not consistently superior. The discussion should
better balance performance and interpretability rather than implying broad performance advantage.
The ablation study is based on a single fold only, which limits the reliability of the
conclusions. The authors should either repeat ablations across multiple folds or clearly tone down
the strength of the claims.
Runtime comparisons need clearer reporting. The manuscript reports runtimes for ablation variants,
but does not compare computational cost against all baseline methods under the same hardware and
settings.
The graphical abstract and Figure~1 contain small text and may be difficult to read. The authors
should improve figure resolution, font size, and visual clarity.
Some figures, especially SHAP visualizations, are too small to interpret in the manuscript.
Larger panels or supplementary high-resolution versions should be provided.
The paper should provide the exact hyperparameters used for LightGBM, CNN, MiniRocket, RDST,
tsfresh, and catch22 to ensure reproducibility.
The authors state that SublimeX is open source, but the manuscript should include a permanent
repository link, version number, and ideally a reproducibility package for all experiments.
The manuscript should better define what ``interpretable'' means in this context. Traceability to
a signal segment is useful, but interpretability also depends on clinical or biological validity,
stability, and expert confirmation.
The SHAP analysis is helpful, but it is only shown for two datasets. The authors should clarify
why REMC and AZT1D were selected and whether similar interpretability patterns hold across other
datasets.
The discussion should include a stronger limitations section. Important limitations include greedy
search, local optima, risk of overfitting, high computational cost, limited baseline comparison,
and reduced generalization in cross-subject settings.
Several writing and formatting issues need correction, including typographical errors such as
``anaylsis,'' inconsistent capitalization of dataset names, and overly long paragraphs in the
Discussion.
The manuscript would benefit from a clearer contribution statement at the end of the
Introduction, listing the main methodological and empirical contributions.
The authors should avoid overstating clinical applicability. The datasets are benchmark datasets,
and no prospective clinical validation or expert user study is presented.
The manuscript should clarify whether SublimeX is intended mainly for prediction, hypothesis
generation, or clinical decision support, because the validation requirements differ for each use
case.
The manuscript is promising and relevant to Pattern Recognition, but it requires major revision
to improve methodological rigor, reproducibility, baseline comparison, and clarity of claims.}

\medskip\noindent\textbf{Our point-by-point responses}

\rcitem{1. Novelty vs.\ interval / shapelet / ROCKET / supervised extractors.}{The manuscript
presents an interesting interpretable feature extraction framework, but the novelty of SublimeX
should be stated more clearly. The authors should explain how it differs substantially from
existing interval-based, shapelet-based, ROCKET/MiniROCKET, and supervised feature extraction
methods.}

\response{We now state the novelty in the Introduction and the new Sec.~1.2 (Related Work),
contrasting SublimeX with fixed libraries, shapelet methods, and deep models (Table~1).}{
  \item Introduction: ``Unlike tsfresh-style fixed libraries, SublimeX searches which regions
    matter for each task.''
  \item Sec.~1.2: shapelet vs.\ segment-mean contrast; ``combines traceable segment means with
    supervised search''; Table~1 compares type, \# features, multivariate handling (joint vs.\
    per-channel), and interpretability (High/Medium/Low) across fixed libraries, interval/shapelet
    methods, ROCKET/MiniRocket, and deep learning baselines.
  \item Methods, Baselines: tsfresh, catch22, MiniRocket, RDST, and 1D CNN under matched protocols.
}

\rcitem{2. Title ``time series and spatial data''.}{The title mentions both ``time series and
spatial data,'' but most of the paper focuses on time-series biomedical datasets. The spatial data
contribution is not sufficiently demonstrated. The authors should either strengthen the spatial-data
evaluation or revise the title and scope.}

\response{We agree that six of seven benchmarks are temporal and that REMC is our sole spatially
indexed evaluation (1D epigenomic profiles along genomic position, not 2D imaging). We kept the
title and state in Methods, Overview that segment search is agnostic to temporal vs.\ spatial
indexing; REMC is evaluated per cell line (56 models, pooled reporting) with full SHAP
interpretation on E003.}{
  \item Methods, Overview: ``The search is agnostic to whether the index axis is temporal or
    spatial: it operates on multivariate 1D sequences, and only the domain interpretation of
    channels and segments differs.''
  \item Methods, Datasets: REMC as ``spatial epigenomic gene-expression prediction \ldots{}
    (56 cell lines)''; 5-fold CV per cell line, scores pooled across cell lines$\times$folds.
  \item Results, Domain-Specific Feature Interpretation: REMC labeled ``spatial epigenomic
    profiles''; SHAP on cell line E003.
  \item Discussion, Limitations: scope limited to 1D genomic locus profiles; no 2D spatial benchmark.
}

\rcitem{3. Abstract competitiveness.}{The abstract is informative, but it overstates
competitiveness. SublimeX is not the best-performing method on several datasets, including MITBIH,
Emotions, MIMIC-IV, PAMAP2, and SVD. The abstract should more cautiously describe the results.}

\response{We toned down competitiveness claims throughout. Overstated wording on AZT1D RMSE,
dataset leadership, and state-of-the-art performance was replaced with cautious language.}{
  \item Abstract, Results (revised): ``ties the best REMC AUC ($0.931 \pm 0.023$)\ldots ranks
    second on AZT1D (RMSE $30.98 \pm 5.65$) and MIMIC-IV (AUC $0.818 \pm 0.024$)''; ``uses
    ${\sim}6$--$20$ descriptors versus thousands for library and kernel methods.''
  \item Results, Performance Comparison (revised): ``MiniRocket achieves the lowest mean test RMSE
    ($30.76{\pm}5.38$); SublimeX is second ($30.98{\pm}5.65$)'' with ``$p{=}0.39$''.
  \item Conclusion (revised): ``Interpretability need not come at the cost of performance when
    feature search is supervised, validation-driven, and constrained to a compact representation''
    (generalized; no per-dataset score recap).
}

\rcitem{4. Stopping criterion and metric direction.}{The methodology needs more detail on the
stopping criterion. The paper states that the loop stops when no candidate improves performance,
but the exact definition of improvement, tolerance threshold, and validation metric direction
should be clearly specified.}

\response{We now define the stopping rule precisely, including strict improvement and metric
direction.}{
  \item Methods, Overview (revised): ``stops once no new segment yields a strict improvement.''
  \item Methods, Iterative Feature Discovery (revised): ``only if it improves the running-best
    validation score''; ``When no candidate yields a strict improvement, the loop terminates.''
  \item Methods, Implementation Details: ``300 trials per proposed feature''; maximizes accuracy or
    AUC, minimizes RMSE on AZT1D.
}

\rcitem{5. Fixed 50/50 inner split.}{The use of a single deterministic 50/50 train/validation
split during feature optimization may introduce instability and overfitting. The authors should
justify this choice and discuss whether nested cross-validation or repeated validation would
provide more reliable feature selection.}

\response{We added a \emph{Resampling} ablation (Table~4) that redraws the fit/selection partition on
every Optuna trial. Resampling was generally slightly worse than the fixed-split baseline across
benchmarks; for example, AZT1D RMSE rose from $30.98 \pm 5.65$ to $31.17 \pm 5.85$ ($p{<}0.05$), and
PAMAP2 LOSO accuracy fell from $0.860 \pm 0.081$ to $0.843 \pm 0.094$. We keep the fixed 50/50
split because it stabilizes candidate scores across trials; Limitations discusses overfitting risk,
and nested validation is future work.}{
  \item Methods, Evaluation: PAMAP2 LOSO; once-assigned 50/50 fit/selection split on training
    subjects; anti-leakage and AZT1D per-patient protocol.
  \item Methods, Iterative Feature Discovery: inner 50/50 split only; refit before test scoring.
  \item Ablation Study (Table~4), \emph{Resampling}: generally slightly below baseline; e.g.\ PAMAP2
    $0.860{\pm}0.081$ vs.\ $0.843{\pm}0.094$; AZT1D RMSE $30.98{\pm}5.65$ vs.\ $31.17{\pm}5.85$.
  \item Discussion, Limitations: fixed split preferred over Resampling; overfitting to optimization
    split, especially PAMAP2.
  \item Discussion, Future Directions: ``nested or repeated internal validation.''
}

\rcitem{6. PAMAP2 internal vs.\ LOSO validation.}{The internal feature-search validation split
does not match the external evaluation design for some datasets, especially PAMAP2 with
leave-one-subject-out validation. This likely explains the large validation-test gap. The authors
should consider using subject-independent validation during feature search.}

\response{We explain the large PAMAP2 validation--test gap (${\sim}26\%$): the fit and selection
halves share the same training subjects, whereas LOSO tests an entirely unseen subject; the
held-out subject never enters the search. \textbf{We did not implement} subject-independent
validation during every search trial, owing to compute and the 35-page scope; the \emph{Resampling}
ablation confirms that redrawn inner splits do not help and hurt PAMAP2 under LOSO.}{
  \item Methods, Evaluation: PAMAP2 LOSO; once-assigned fit/selection split on training subjects;
    held-out subject never enters search.
  \item Methods, Evaluation / Results: inner split not subject-independent; validation--test gap
    reported in Results; interpretation in Discussion.
  \item Results, Incremental Feature Analysis: PAMAP2 gap ${\sim}26\%$ (numbers only).
  \item Discussion, Summary of Main Findings: PAMAP2 fit/selection vs.\ LOSO interpretation.
  \item Table~4: eight-fold PAMAP2 LOSO ablations with mean$\pm$SD; \emph{Resampling} underperforms
    baseline.
  \item Discussion, Limitations: fixed inner split preferred over \emph{Resampling}; subject-holdout
    validation during every search trial omitted for compute cost.
}

\rcitem{7. Leakage / outer-fold-only search.}{The manuscript should clarify whether feature
extraction is performed independently within each training fold only. If any information from
validation or test folds influences feature discovery, this would create data leakage.}

\response{Methods now state explicitly that feature discovery uses the training partition only.
Test folds, LOSO subjects, and AZT1D temporal test windows never enter segment search.}{
  \item Methods, Evaluation (new): ``never enters segment search''; ``no leakage from test data into
    feature discovery''; 5-fold CV, PAMAP2 LOSO, AZT1D temporal split.
  \item Methods, Iterative Feature Discovery: fit/selection split only (cross-reference to Evaluation).
}

\rcitem{8. HIVE-COTE 2.0, InceptionTime, LSTM/GRU.}{The comparison with baselines is incomplete.
Important state-of-the-art time series methods such as HIVE-COTEv2, InceptionTime, and
LSTM/GRU-based models are only mentioned as limitations but not evaluated. At minimum, the authors
should justify why these were excluded.}

\response{We justify the baseline scope in Methods and Limitations: our focus is interpretable,
segment-traceable features under matched biomedical splits, so we compare against five complementary
baselines (including a 1D CNN on the same folds) and explain why HIVE-COTE, InceptionTime, LSTM, and
the UCR archive were excluded.}{
  \item Methods, Baselines: ``interpretable, segment-traceable features''; hybrid/deep/recurrent
    methods cited as not directly comparable; 1D CNN on same folds.
  \item Discussion, Limitations: HIVE-COTEv2, InceptionTime, LSTM ``not included''; UCR archive
    ``do not evaluate on.''
}

\rcitem{9. Statistical significance.}{The authors should provide statistical significance testing
across methods. Many reported differences are very small, such as AZT1D RMSE 31.11 vs.\ 31.18 and
REMC AUC 0.931 vs.\ 0.931, so claims of superior performance require statistical support.}

\response{We added paired Wilcoxon tests and re-ran AZT1D reporting. Table~3 now gives
mean$\pm$SD of per-patient \emph{test} RMSE rather than a single pooled summary.
MiniRocket is best ($30.76 \pm 5.38$); SublimeX is second ($30.98 \pm 5.65$, $p{=}0.39$ vs.\
MiniRocket).}{
  \item Methods, Evaluation: ``paired two-sided Wilcoxon tests ($\alpha{=}0.05$)''; ``AZT1D
    per-patient test RMSE.''
  \item Table~3 (revised): asterisks vs.\ dataset-best method; AZT1D reported as mean$\pm$SD across
    23 patients.
  \item Results, Performance Comparison (revised): ``does not differ significantly from MiniRocket
    ($p{=}0.39$).''
}

\rcitem{10. Balance performance vs.\ interpretability.}{The results show that SublimeX is
competitive but not consistently superior. The discussion should better balance performance and
interpretability rather than implying broad performance advantage.}

\response{The Discussion now balances performance against interpretability, using measured language
(MiniRocket best on AZT1D, SublimeX second without a significant gap) and framing interpretability
as segment traceability on public benchmarks rather than a performance advantage.}{
  \item Discussion, Summary of Main Findings: opens ``In this study we presented \ldots{}
    SublimeX''; ``competitive'' / ``stays competitive'' wording; MiniRocket best on AZT1D,
    SublimeX second; baselines lead on several other tasks; PAMAP2 LOSO gap interpreted here
    (not in Incremental Feature Analysis).
  \item Abstract, Conclusions: ``Interpretability need not sacrifice performance'' with traceable
    channel/transform/segment descriptors.
  \item Discussion, Comparison to Related Work: predictions related to measurable signal structure;
    segment descriptors for inspection/SHAP (no deployment or intended-use claim).
  \item Discussion, Limitations: performance--interpretability trade-offs, baseline scope,
    cross-subject generalization (PAMAP2).
}

\rcitem{11. Multi-fold ablations.}{The ablation study is based on a single fold only, which limits
the reliability of the conclusions. The authors should either repeat ablations across multiple
folds or clearly tone down the strength of the claims.}

\response{We extended ablations from a single fold to all folds, as reviewers requested.
The revision reports mean$\pm$SD with Wilcoxon markers in Table~4. REMC ablations use cell line
E003 only and are not comparable to the pooled 56-line SublimeX score in Table~3. }{
  \item Ablation Study (revised): ``report mean$\pm$standard deviation across folds''; REMC E003
    only; Wilcoxon asterisks in Table~4; \emph{Resampling} variant added.
  \item Table~4: multi-fold scores, feature counts, and runtimes for eight variants (including
    \emph{Resampling}).
}

\rcitem{12. Runtime vs.\ all baselines.}{Runtime comparisons need clearer reporting. The
manuscript reports runtimes for ablation variants, but does not compare computational cost against
all baseline methods under the same hardware and settings.}

\response{Table~3 now reports end-to-end wall-clock time (feature extraction plus training) for
SublimeX and all five baselines, measured on the single workstation stated in Implementation
Details.}{
  \item Table~3: wall-clock time row for SublimeX and all five baselines.
  \item Methods, Implementation Details: ``Experiments were run on a MacBook Pro M2 (12 cores,
    32\,GB RAM).''
  \item Results, Performance Comparison: ``SublimeX is the slowest method'' on large datasets;
    modest runtime on shorter series.
  \item Discussion, Limitations: slower than several baselines despite fewer features.
}

\rcitem{13--14. Graphical abstract and Figure~1.}{The graphical abstract and Figure~1 contain
small text and may be difficult to read. The authors should improve figure resolution, font size,
and visual clarity.}

\response{We redrew Figure~1 with larger type and a simplified loop.}{
  \item Figure~1: full-width methodology flowchart with enlarged lettering.
  \item Methods, Overview: ``stops once no new segment yields a strict improvement.''
}

\rcitem{15. SHAP figure size.}{Some figures, especially SHAP visualizations, are too small to
interpret in the manuscript. Larger panels or supplementary high-resolution versions should be
provided.}

\response{We enlarged Figure~4 to a full-width, six-panel SHAP layout for REMC and AZT1D so panels
are legible at print scale. }{
  \item Results, Domain-Specific Feature Interpretation: ``full six-panel SHAP layouts for REMC
    (spatial epigenomic profiles) and AZT1D (glucose forecasting) 
    with larger font sizes and only top five features for SHAP beeswarm plot''
}

\rcitem{16. Hyperparameters.}{The paper should provide the exact hyperparameters used for
LightGBM, CNN, MiniRocket, RDST, tsfresh, and catch22 to ensure reproducibility.}

\response{We now list all hyperparameters explicitly: LightGBM (depth~5, shared readout) in
Implementation Details; tsfresh, catch22, MiniRocket, RDST, and CNN settings in Methods,
Baselines.}{
  \item Methods, Baselines: tsfresh, catch22, MiniRocket, RDST, and CNN settings in prose
    (replacing ``default hyperparameters'' and the former supplementary hyperparameter table).
}

\rcitem{17. Open source and reproducibility.}{The authors state that SublimeX is open source, but
the manuscript should include a permanent repository link, version number, and ideally a
reproducibility package for all experiments.}

\response{We now state version, installation, and repository URLs explicitly in Methods. The
linked manuscript repository documents end-to-end reproduction of benchmarks, ablations, tables,
and figures. Dataset preprocessing is in Supplementary Data.}{
  \item Methods, Implementation Details (revised): ``published as an open-source Python package on
    PyPI (version~0.1.2)''; installable via pip; manuscript GitHub URL for reproduction. 
    We ensure the permannence of the link and maintain the Python package as an actively developed package.
  \item Manuscript repository \url{https://github.com/Prgrmmrjns/SublimeX-manuscript}: 
  instructions to regenerate tables/figures from saved features or rerun the full experiment pipeline.
  Note the repo does not contain any datasets, these have to be downloaded beforehand. 
  We provide instructions in the repo and in the Supplementary Data for downloading and reproducing the data
  as well as the corresponding preprocessed files.
}

\rcitem{18. Definition of ``interpretable''.}{The manuscript should better define what
``interpretable'' means in this context. Traceability to a signal segment is useful, but
interpretability also depends on clinical or biological validity, stability, and expert
confirmation.}

\response{We now define interpretability operationally in the manuscript: a feature is
interpretable when it is \emph{traceable}---each descriptor names its channel, transform, segment,
and aggregation---and when that choice is \emph{stable across folds} and \emph{plausible relative to
known biology} where benchmarks allow. We do not claim formal expert confirmation or clinical
readiness; SHAP on REMC and AZT1D is our check that ranked segments align with established domain
mechanisms.}{
  \item Introduction: compact features with ``explicit descriptions that experts can inspect
    directly''; objective~(iii) assesses interpretability via post-hoc analyses.
  \item Results, Transform Composition: fold-stable transform mix and first-feature segment
    positions (Figure~2).
  \item Results, Domain-Specific Feature Interpretation: ``Each SublimeX feature is readable as a
    named channel, transform, segment, and aggregation, and we check interpretability by asking
    whether those choices stay stable across folds and line up with SHAP rankings and known
    biology where benchmarks allow''; REMC epigenomic hierarchy; AZT1D mean reversion then
    insulin/carbohydrate inputs (Figure~4).
  \item Discussion: ``traceable, and thereby interpretable''; SHAP recovers meaningful hierarchies
    ``without any domain guidance.''
}

\rcitem{19. SHAP on two datasets only.}{The SHAP analysis is helpful, but it is only shown for
two datasets. The authors should clarify why REMC and AZT1D were selected and whether similar
interpretability patterns hold across other datasets.}

\response{We chose REMC and AZT1D for full SHAP figures because they have the clearest domain
ground truth and represent the two modalities in the title. For the remaining benchmarks we provide
transform-stability (Figure~2) and incremental-feature analyses (Figure~3). 
There was no space to add SHAP on the remaining datasets. 
We hope that with our in-depth walkthrough for these two datasets the reader can appreciate the advantages of segment-based traceable features
in combination with SHAP analysis for domain interpretation.}{
  \item Results, Domain-Specific Feature Interpretation: expanded REMC and AZT1D subsections
    (biology/clinical context, panel-by-panel SHAP walkthrough); AZT1D mean reversion narrative.
}

\rcitem{20. Stronger limitations section.}{The discussion should include a stronger limitations
section. Important limitations include greedy search, local optima, risk of overfitting, high
computational cost, limited baseline comparison, and reduced generalization in cross-subject
settings.}

\response{The Limitations are now a dedicated, multi-paragraph subsection, with Future Directions
kept separate.}{
  \item Discussion, Limitations: greedy search; runtime; PAMAP2 gap and \emph{Resampling}
    ablation; HIVE-COTE/InceptionTime/LSTM; UCR; slower than lighter baselines; benchmark scope
    (no UCR).
  \item Discussion, Future Directions: nested validation; learned transforms; efficient search.
}

\rcitem{21. Writing and formatting.}{Several writing and formatting issues need correction,
including typographical errors such as ``anaylsis,'' inconsistent capitalization of dataset names,
and overly long paragraphs in the Discussion.}

\response{We corrected the typographical errors, standardized dataset names, moderated performance
claims, and tightened the wording and flow throughout.}{
  \item Introduction and Methods: Sec.~1.1 Objectives + Sec.~1.2 Related Work; smoother
    transitions and reduced repetition.
  \item Discussion: shorter paragraphs; PAMAP2 interpretation in Summary; Limitations and Future
    Directions separated.
  \item Conclusion: three generalized sentences (no dataset-specific score recap).
  \item Abstract and body: ``MIT-BIH'' (not ``MITBIH'').
}

\rcitem{22. Contribution statement.}{The manuscript would benefit from a clearer contribution
statement at the end of the Introduction, listing the main methodological and empirical
contributions.}

\response{We added an explicit contribution statement: Sec.~1.1 (Objectives) lists three numbered
aims, followed by Sec.~1.2 (Related Work), both before Methods.}{
  \item Sec.~1.1: ``The objectives of our study are as follows: (i) \ldots{} (ii) \ldots{}
    (iii) \ldots{}''
  \item Sec.~1.2: Related Work and Table~1 (method taxonomy, multivariate handling, interpretability)
    before Methods.
}

\rcitem{23. Clinical applicability.}{The authors should avoid overstating clinical applicability.
The datasets are benchmark datasets, and no prospective clinical validation or expert user study
is presented. The manuscript should clarify whether SublimeX is intended mainly for prediction,
hypothesis generation, or clinical decision support, because the validation requirements differ
for each use case.}

\response{We do not assign SublimeX to prediction, hypothesis generation, or clinical decision
support, and we do not claim clinical readiness. We state the need for inspectable predictions
indirectly: many biomedical tasks require outputs tied to measurable signal structure, and SublimeX
yields segment-traceable descriptors (channel, transform, location) that support such inspection on
our public benchmarks---illustrated with SHAP and fold stability, not prospective validation.}{
  \item Abstract, Introduction (revised): ``inform clinical and scientific analysis'' (motivation
    only; not a claim about intended use).
  \item Discussion, Comparison to Related Work (revised): removed ``Biomedical deployment places
    constraints'' and ``attractive for biomedical use''; now ``Many biomedical modeling tasks require
    predictions that can be related to measurable signal structure \ldots{} SublimeX represents each
    sample through segment-based features with explicit, human-readable descriptions.''
  \item Discussion, Comparison to Related Work: MIMIC sentence now ``segment-traceable features tied
    to identifiable ICU signal regions'' (replacing ``support clinical validation'').
}

\rcitem{Reviewer \#2 --- metareview}{The manuscript presents an interesting and relevant approach
to interpretable feature extraction for time series and spatial data, addressing an important gap
between predictive performance and model transparency. The proposed framework is well motivated and
demonstrates competitive results across multiple datasets, with a clear emphasis on
interpretability. However, several aspects require improvement to strengthen the contribution. In
particular, the objectives and novelty of the method should be more explicitly stated, and the
distinction from existing approaches clarified. The experimental section would benefit from more
rigorous statistical analysis, clearer reporting of reproducibility details, and stronger
justification of evaluation protocols. Additionally, some claims regarding performance and
interpretability should be moderated or better supported by quantitative evidence. The manuscript
is generally well structured, but certain sections could be streamlined for clarity, and minor
language editing is recommended. Overall, the work has merit, but revisions are needed to improve
clarity, rigor, and positioning.}

\response{We thank Reviewer~2 for this constructive summary. Every theme is addressed in Q1--Q8 and
in the point-by-point responses below, with quoted wording from the revised manuscript. Where a
request could not be fully met, we explain the constraint (chiefly the 35-page limit) and what we
provided instead.}{
  \item See Q1--Q8 and Reviewer~1 items 1--23.
}

\vspace{1em}
\noindent We hope this letter and the revised manuscript address the reviewers' concerns.

\vspace{0.8em}
\noindent Sincerely,\\
The Authors

\end{document}
